\chapter{On-Board Software Architecture \& Autonomy Stack}
\label{ch:software_autonomy}

\section{System-Level Software Topology (PX4 meets ROS 2)}
\label{sec:software_topology}
% Erklärung der zwei Welten: PX4 auf dem Pixhawk (Echtzeit/uORB) 
% und ROS 2 Humble auf dem Jetson Orin Nano (Linux/Missionslogik).

\section{Inter-Processor Bridge (micro-XRCE-DDS)}
\label{sec:xrce_dds_bridge}
% Beschreibung der Datenbrücke über Ethernet zwischen Autopilot und Companion Computer.

\section{ROS 2 Node Architecture \& Communication Matrix}
\label{sec:ros2_node_architecture}
% Detaillierte Beschreibung aller Nodes (ai_perception, target_geolocator, swarm_coordinator)
% inkl. der Tabelle mit Publishers, Subscribers und Topics.

\section{AI Perception Pipeline (YOLO \& TensorRT)}
\label{sec:ai_perception_pipeline}
% Einbindung des EO/IR-RTSP-Videostreams, Hardware-Beschleunigung auf dem Jetson
% und Transformation von Pixel-Koordinaten in reale GPS-Daten (WGS84/MGRS).

\section{Decentralized Swarm Logic \& Multi-Agent Framework}
\label{sec:swarm_logic}
% Wie die Drohnen dezentral via Namespaces (/uav_N/) und P2P-Protokoll 
% das Suchgebiet aufteilen und Personen-Fundmeldungen austauschen.

\section{Multi-Vehicle SITL Simulation (Gazebo \& PX4)}
\label{sec:gazebo_multi_vehicle_sim}
% Nachweis der Funktionsfähigkeit der Schwarm-Software mit 3 simulierten VTOLs in Gazebo.